\section{条件概率}

\subsection{条件概率}

\begin{definition}
    对于事件 $A, B$，若 $P(A) > 0$，定义
    $$
    P(B \mid A) = \frac{P(A B)}{P(A)}
    $$
    为在事件 $A$ 发生的条件下事件 $B$ 发生的概率。
\end{definition}

条件概率 $P(\cdot \mid A)$ 满足下面的性质：

\begin{itemize}
    \item 非负性：对于任意事件 $B$ 都有 $P(B \mid A) \ge 0$；
    \item 规范性：对于必然事件 $S$ 有 $P(S \mid A) = 1$；
    \item 可列可加性：设 $B_1, B_2,\dots$ 是两两互不相容的事件，那么
    $$
    P\ab(\bigcup_{i = 1}^{\infty} B_i \mid A) = \sum_{i = 1}^{\infty} P(B_i \mid A)
    $$
\end{itemize}

将条件概率进行变形，可以得到\textbf{乘法公式}：

\begin{theorem}[乘法定理]
    设 $P(A) > 0$，那么有：
    $$
    P(A B) = P(B \mid A) P(A)
    $$
\end{theorem}

\subsection{Bayes 公式}

先定义样本空间的划分：

\begin{definition}[划分]
    设 $S$ 为试验 $E$ 的样本空间，若一组事件 $B_1, B_2,\dots, B_n$ 满足：
    \begin{itemize}
        \item $\forall\,i \neq j: B_i \cap B_j = \varnothing$；
        \item $\displaystyle \bigcup_{i = 1}^{n} B_i = S$。
    \end{itemize}
    那么称 $B_1, B_2,\dots, B_n$ 是对样本空间 $S$ 的\textbf{划分}。
\end{definition}

然后我们有全概率公式：

\begin{theorem}[全概率公式]
    设 $A$ 是事件，$B_1, B_2,\dots, B_n$ 是对样本空间 $S$ 的划分，那么：
    $$
    P(A) = \sum_{i = 1}^{n} P(A \mid B_i) P(B_i)
    $$
    \begin{proof}
        $$
        \begin{aligned}
            P(A) & = P(A S) = P\ab(A \bigcup_{i = 1}^{n} B_i) \\
            & = P\ab(\bigcup_{i = 1}^{n} A B_i) \\
            & = \sum_{i = 1}^{n} P(A B_i) = \sum_{i = 1}^{n} P(A \mid B_i) P(B_i)
        \end{aligned}
        $$
    \end{proof}
\end{theorem}

根据全概率公式进行变形，可以得到 Bayes 公式：

\begin{theorem}
    设 $A$ 是事件，$B_1, B_2,\dots, B_n$ 是对样本空间 $S$ 的划分，那么：
    $$
    P(B_i \mid A) = \frac{P(A \mid B_i) P(B_i)}{\sum_{j = 1}^{n} P(A \mid B_j) P(B_j)}
    $$
\end{theorem}

\subsection{独立性}

\begin{theorem}[独立性]
    设 $A, B$ 是两个事件，若
    $$
    P(A B) = P(A) P(B)
    $$
    那么称事件 $A, B$ \textbf{相互独立}。
\end{theorem}

当两个事件相互独立时，一个事件是否发生不会影响另一个事件发生的概率，即
$$
P(B \mid A) = P(B \mid \overline{A}) = P(B)
$$

\subsection{作业}

\begin{problem}
    习题 14
    \begin{solution}
        \begin{enumerate}
            \item[\textbf{1)}]
            $$
            \begin{aligned}
                \origin & = \frac{P(B (A \cup \overline{B}))}{P(A \cup \overline{B})} = \frac{P(A B)}{P(A \cup \overline{B})} \\
                & = \frac{P(A) - P(A \overline{B})}{P(A) + P(\overline{B}) - P(A \overline{B})} = \frac{1}{4}
            \end{aligned}
            $$
            \item[\textbf{2)}]
            $$
            \begin{aligned}
                \origin & = P(A) + P(B) - P(A B) \\
                & = P(A) + \frac{P(A B)}{P(A \mid B)} - P(A B) \\
                & = P(A) + P(A) P(B \mid A) \ab(\frac{1}{P(A \mid B)} - 1) \\
                & = \frac{1}{3}
            \end{aligned}
            $$
        \end{enumerate}
    \end{solution}
\end{problem}

\begin{problem}
    习题 17
    \begin{solution}
        样本空间大小为 $10 \cdot 9 = 90$，于是：
        \begin{enumerate}
            \item[\textbf{1)}] $P = \frac{8 \cdot 7}{90} = \frac{28}{45}$；
            \item[\textbf{2)}] $P = \frac{2 \cdot 1}{90} = \frac{1}{45}$；
            \item[\textbf{3)}] $P = \frac{8 \cdot 2 + 2 \cdot 8}{90} = \frac{16}{45}$；
            \item[\textbf{4)}] $P = \frac{8 \cdot 2 + 2 \cdot 1}{90} = \frac{1}{5}$。
        \end{enumerate}
    \end{solution}
\end{problem}

\begin{problem}
    习题 18
    \begin{solution}
        每次拨号正确的概率为 $\frac{1}{10}$，因此不超过三次拨号正确的概率为：
        $$
        P = \frac{1}{10} + \frac{9}{10} \cdot \frac{1}{10} + \frac{9}{10} \cdot \frac{9}{10} \cdot \frac{1}{10} = \frac{271}{1000}
        $$
        若已知最后一个数字是奇数，那么每次拨号正确概率为 $\frac{1}{5}$，不超过三次拨号正确概率为：
        $$
        P' = \frac{1}{5} + \frac{4}{5} \cdot \frac{1}{5} + \frac{4}{5} \cdot \frac{4}{5} \cdot \frac{1}{5} = \frac{61}{125}
        $$
    \end{solution}
\end{problem}

\begin{problem}
    习题 21
    \begin{solution}
        设事件 $A$ 表示『此人为男性』，事件 $B$ 表示『此人为色盲患者』，那么可知：
        $$
        P(A) = P(\overline{A}) = \frac{1}{2},\ P(B \mid A) = \frac{1}{20},\ P(B \mid \overline{A}) = \frac{1}{400}
        $$
        题目所求概率为
        $$
        P(A \mid B) = \frac{P(B \mid A) P(A)}{P(B \mid A) P(A) + P(B \mid \overline{A}) P(\overline{A})} = \frac{20}{21}
        $$
    \end{solution}
\end{problem}

\begin{problem}
    习题 22
    \begin{solution}
        记事件 $A$ 表示『第一次及格』，事件 $B$ 表示『第二次及格』。

        \begin{enumerate}
            \item[\textbf{1)}]
            $$
            P = P(A \cup B) = P(A) + P(B \mid \overline{A}) P(\overline{A}) = \frac{3}{2} p - \frac{p^2}{2}
            $$
            \item[\textbf{2)}]
            $$
            P(A \mid B) = \frac{P(B \mid A) P(A)}{P(B \mid A) P(A) + P(B \mid \overline{A}) P(\overline{A})} = \frac{2 p^2}{p^2 + p}
            $$
        \end{enumerate}
    \end{solution}
\end{problem}

\begin{problem}
    习题 22
    \begin{solution}
        记事件 $A$ 表示『选中了第一箱』，事件 $B$ 表示『第一次取到的零件是一等品』，事件 $C$ 表示『第二次取到的零件是一等品』。
        \begin{enumerate}
            \item[\textbf{1)}]
            $$
            P(B) = P(B \mid A) P(A) + P(B \mid \overline{A}) P(\overline{A}) = \frac{1}{5} \cdot \frac{1}{2} + \frac{3}{5} \cdot \frac{1}{2} = \frac{2}{5}
            $$
            \item[\textbf{2)}]
            $$
            P(C \mid B) = \frac{P(B C)}{P(B)} = \frac{P(A B C) + P(\overline{A} B C)}{P(A B) + P(\overline{A} B)} = \frac{\frac{1}{2} \cdot \frac{1}{5} \cdot \frac{9}{49} + \frac{1}{2} \cdot \frac{3}{5} \cdot \frac{17}{29}}{\frac{1}{2} \cdot \frac{1}{5} + \frac{1}{2} \cdot \frac{3}{5}} = \frac{690}{1421}
            $$
        \end{enumerate}
    \end{solution}
\end{problem}

\begin{problem}
    习题 28
    \begin{solution}
        \begin{enumerate}
            \item[\textbf{1)}] $P_1 = 0.8 \cdot 0.9 = 0.72$；
            \item[\textbf{2)}] $P_2 = 1 - (1 - 0.8) (1 - 0.9) = 0.98$；
            \item[\textbf{3)}] $P_3 = 0.8 \cdot (1 - 0.9) + (1 - 0.8) \cdot 0.9 = 0.26$；
        \end{enumerate}
    \end{solution}
\end{problem}

\begin{problem}
    习题 31
    \begin{solution}
        \begin{enumerate}
            \item[\textbf{1)}] 可能对。若$A, B$ 互不相容那么 $P(A B) = 0$，若 $A, B$ 独立那么 $P(A B) = P(A) P(B)$，只有 $A = S \lor A = \varnothing$ 时是对的；
            \item[\textbf{2)}] 可能对。$A = S$ 或 $A = \varnothing$。
            \item[\textbf{3)}] 必然错。因为若 $A, B$ 不相容那么 $P(A) + P(B) \le 1$。
            \item[\textbf{4)}] 可能对。 
        \end{enumerate}
    \end{solution}
\end{problem}

\begin{problem}
    习题 34
    \begin{solution}
        \begin{enumerate}
            \item[\textbf{1)}]
            $$
            R_1 = p_1 \ab(1 - (1 - p_2 p_3) (1 - p_4)) = p_1 (p_2 p_3 + p_4 - p_2 p_3 p_4)
            $$
            \item[\textbf{2)}] 当已知元件 3 失效时：
            $$
            R' = 1 - \ab(1 - p^2)^2 = 2 p^2 - p^4
            $$
            当已知元件 3 正常时，只要 1, 4 其中一个正常、2, 5 其中一个正常即可：
            $$
            R'' = \ab(1 - \ab(1 - p^2))^2 = \ab(2 p - p^2)^2
            $$
            根据全概率公式：
            $$
            R_2 = (1 - p) R' + p R'' = 2 p^2 + 2 p^3 - 5 p^4 + 2 p^5
            $$
        \end{enumerate}
    \end{solution}
\end{problem}

\begin{problem}
    习题 40
    \begin{solution}
        记事件 $A$ 为『输入是 \texttt{AAAA}』，事件 $B$ 为『输入是 \texttt{BBBB}』，事件 $C$ 为『输入是 \texttt{CCCC}』，事件 $X$ 为『输出是 \texttt{ABCA}』。那么可知：
        $$
        \begin{aligned}
            P(X \mid A) = \alpha^2 \ab(\frac{1 - \alpha}{2})^2 \\
            P(X \mid B) = \alpha \ab(\frac{1 - \alpha}{2})^3 \\
            P(X \mid C) = \alpha \ab(\frac{1 - \alpha}{2})^3 \\
        \end{aligned}
        $$
        根据 Bayes 公式可得所求概率为：
        $$
        P(A \mid X) = \frac{P(X \mid A) P(A)}{P(X \mid A) P(A) + P(X \mid B) P(B) + P(X \mid C) P(C)} = \frac{2 \alpha p_1}{2 \alpha p_1 + (1 - p_1) (1 - \alpha)}
        $$
    \end{solution}
\end{problem}